\section{Puissance}
\subsection{Puissance}
\begin{definitionbox}
La puissance d’un nombre relatif est le produit de plusieurs facteurs égaux : 
$$ a^n = \underbrace{a\times a \times ... \times a}_{\text{\normalfont{n-fois}}}$$
ou $a$ est appelé la \bb{base} et $n$ est appelé \bb{l'exposant}.
\end{definitionbox}

\begin{exemplebox}
    \begin{itemize}
        \item $3^4 = 3 \times 3\times 3\times 3 = 81$
        \item $(-2)^3 = (-2) \times (-2)\times (-2) = -8$
    \end{itemize}
\end{exemplebox}

\begin{exercicebox}
    \begin{itemize}
        \item $(-8)^2 = \answer{(-8) \times (-8) = 64}$
        \item $5^3 = \answer{5 \times 5\times 5 = 125}$
    \end{itemize}
\end{exercicebox}

\subsection{Règle des signes}
\begin{propbox}
\bb{Signe de la puissance d'un nombre}
\begin{enumerate}
    \item la puissance d'un nombre positif est \bb{toujours positive}
    \item la puissance d'un nombre négatif est :
    \begin{itemize}
        \item positive si l'\bb{exposant} est \bb{pair}.
        \item négative si l'\bb{exposant} est \bb{impair}.
    \end{itemize}
\end{enumerate}
\end{propbox}
